
\begin{center}
{\huge \textbf{Journée type}}\\ \vspace{0.5em}

\renewcommand{\arraystretch}{1.40}
\begin{tabular}{|C{3em}|C{12em}|}
\hline
\textbf{Heure} & \textbf{Activité} \\
\hline
7h30  & Réveil, petit déjeuner et mise en condition de navigation des bateaux \\
\hline
8h30  & Brief de navigation + conseil de flotille \\
\hline
9h00  & Appareillage, encas en navigation \\
\hline
12h30 & Déjeuner et repos \\
\hline
14h00 & Deuxième partie de navigation, goûter en mer \\
\hline
17h00 & Arrivée au port, amarage \\
\hline
17h30 & Débriefing de navigation \\
\hline
18h00 & Douche et préparation du repas \\
\hline
19h00 & Repas \\
\hline
20h30 & Veillée. En cas de mouillage, la veillée est avancée à 20h et la vaisselle est faite au retour de la veillée. \\
\hline
22h00 & Chaque jeune est dans son bateau \\
\hline
22h30 & Silence sur le camp \\
\hline
\end{tabular}
\end{center}

\vspace{1cm}
La journée type dimensionne la durée des navigations. 
Elle a néanmoins vocation à être ajustée lors des nuits au mouillage, permettant ainsi de descendre à terre pour les veillées, tout en étant de retour aux bateaux avant le coucher du soleil.
Les équipages de navigation sont autant que possible mixtes afin de varier les découvertes et d'homogénéiser les niveaux de navigation. 
Les deux navigations de sortie de mouillage se feront toutefois en équipe de vie. 
